% !TeX root = ../main.tex

\chapter{绪论}
\section{研究背景与意义}
随着经济社会的快速发展，城市化进程不断加快，交通运输需求日益增加。

交通拥堵成为当前大多数城市面临的一个普遍问题，尤其是大城市，超负荷的交通流量和有限的路网，造成了高峰时段的严重拥堵。根据国内外的研究，交通拥堵不仅仅降低人们的出行效率，还带来了一系列社会经济问题。例如，在北京等一线城市，通勤高峰期的平均车程可能是平时的2-3倍 \parencite{2021_贾若_交通拥堵判别方法研究综述}，导致上班族每天浪费在交通上的时间达到数小时。据估算，全球每年因交通拥堵造成的经济损失已达到数万亿美元。长时间的交通拥堵不仅不仅造成了巨大的能源浪费，也增加了汽车尾气排放，加剧了环境污染。研究揭示，交通拥堵现象每年释放出巨量的二氧化碳及其他有害气体，这不仅对全球气候变化产生了深远影响，也给城市空气质量带来了沉重负担。此外，交通拥堵还直接导致居民生活质量的显著下降，并伴随着交通事故的频繁发生，从而对公共安全构成了更为严峻的挑战。更为严重的是，在高峰时段，交通拥堵限制了公共交通系统的有效承载能力，促使人们更加倾向于依赖私家车出行，这一行为反过来又进一步加剧了交通压力，形成恶性循环。因此，如何准确预测交通流量、优化交通管理、缓解交通拥堵、提高道路利用率，成为了各大城市交通管理部门和研究机构亟待解决的问题。

2019 年中共中央、国务院印发了《交通强国建设纲要》，纲要指出强化前沿关键科技研发，大力发展智慧交通。推动大数据、互联网、人工智能和超级计算等技术与交通行业的深度融合 \parencite{2019_中共中央国务院印发_交通强国建设纲要}。为了提高交通管理效率，许多城市开始借助交通大数据技术进行智能交通管理。交通大数据技术通过收集来自路面传感器、摄像头、GPS等设备的数据，分析交通流量变化，预测未来的交通趋势，从而实现对交通的智能调度和管理。交通流量预测是交通大数据技术中的一个重要应用领域。通过大数据分析，能够对历史交通流量数据进行建模，结合实时数据，提供准确的交通流量预测结果。交通流量预测研究不仅具有重要的理论价值，也对解决城市交通拥堵、提高交通管理水平、促进可持续发展具有重要的现实意义。通过采用大数据和人工智能等技术，城市交通管理可以实现更高效的资源配置和更智能的出行方案。未来，随着技术的不断进步，交通流量预测将在缓解交通压力、减少能源消耗、提高生活质量等方面发挥越来越重要的作用。对于交通流量预测的深入研究，将为城市交通管理提供强有力的决策支持，推动智能交通的普及与发展，最终实现智慧城市的愿景。

交通预测的难点主要在于时空动态变化，交通流量受多种因素的影响，具有高度的随机性和非线性，这使得对交通状态的预测变得异常复杂。为了解决这一问题，许多研究者在交通预测领域进行了广泛的探索，提出了各种方法。传统的统计方法和机器学习方法在处理交通流量数据时，尽管能够在一定程度上反映交通的变化趋势，但仍然存在显著的不足，模型预测效果很大程度上取决于参数预设且对历史数据要求较高，尤其是在应对交通流量这种复杂的随机非线性数据时，预测精度和鲁棒性都受到限制。浅层神经网络虽然能对数据进行拟合，但难以有效捕捉到交通流中的深层次规律。
近年来，基于交通参与对象特点的现代方法逐渐引入了图神经网络等技术。这些方法利用图结构来建模交通路网，能够更好地反映道路之间的空间依赖关系和交通流的时空动态。然而，现有的图方法仍然面临一些挑战，特别是在忽略路网节点间的动态交互效应、序列学习时的梯度爆炸/消失问题以及难以有效学习交通流的非线性扩散机制等方面。为了弥补这些不足，本文提出了一种结合注意力机制与图神经网络的模型，以期通过对节点的动态特征进行加权学习，从而提升交通流量预测的准确性。此外，本文还提出了一种自适应图模型，旨在捕捉交通节点间的潜在关系，为解决现有模型的不足提供新思路。

\section{研究现状}

交通预测方法的分类在现有文献中呈现两种主要形态：动态仿真与数据驱动型方法。基于数学工具（微分方程等）和物理原理的运用，交通流量预测得以通过计算模拟实现的是动态仿真方法\parencite{2015_Vlahogianni_Computational_intelligence_and_optimization_for_transportation_big_data_challenges_and_opportunities}。系统编程的复杂性伴随着该方法的实施过程，计算能力的巨大消耗亦不可忽视，稳态的达成由此受到制约。近年来显著进步的交通数据采集与存储技术使得研究者的关注点发生了转移，数据驱动建模方法正逐渐成为主流选择方向。
\subsection{经典统计模型的预测方法}

鉴于交通流量数据的时序特征，早期研究中多采用时间序列方法对交通流量数据进行统计建模，再基于模型对未来交通流量进行预测。\textcite{2003_Williams_Modeling_and_forecasting_vehicular_traffic_flow_as_a_seasonal_ARIMA_process_Theoretical_basis_and_empirical_results} 基于交通数据非平稳性，提出使用ARIMA方法对数据进行建模和预测. 鉴于ARIMA只能描述单个时间序列的变换，无法同时多个相关时序的动态演变过程，而VAR（向量自回归） 模型不仅能够捕捉单个变量的时间依赖性，还能够揭示多个变量之间的反馈效应 \parencite{1980_Sims_Macroeconomics_and_Reality}。例如，在交通流量预测中，不同道路的流量可能会相互影响，VAR 模型可以帮助分析这些道路流量之间的相互依赖，并进行联合预测, 如\textcite{2022_孙天瑀_基于向量自回归的城市交通流量预测算法研究} 构建交通网络VAR模型实现多节点联动预测，通过特定的向量自回归过程, 可以更好地识别出不同节点之间的依赖模式，有效提升了交通预测的准确性。卡尔曼滤波由Kalman在1960年提出的模型 \parencite{1960_Kalman_A_New_Approach_to_Linear_Filtering_and_Prediction_Problems}，是一种广泛应用于动态系统状态估计的递归算法，特别适用于噪声较大的时间序列预测。该方法基于贝叶斯理论, 通过对系统的状态进行连续的预测和更新，能够有效估计系统的真实状态，并减少测量噪声的影响。\textcite{2021_张塽旖_噪声免疫CatBoost与卡尔曼滤波短期交通流预测研究} 将卡尔曼滤波与CatBoost融合, 并应用于交通流量预测，融合模型能够分辨哪些是噪声, 哪些是反映交通流突变的细微线索。

以上几种模型或方法的研究，都是以数理统计为基础的。这一类模型有些需要严谨的超参数设置，在VAR模型中，每个变量都与自身和其他变量的滞后值进行回归。因此，模型的复杂性取决于所选的滞后期数。滞后期数的选择通常依赖于数据的特征和经验判断，一般通过信息准则（如AIC、BIC）来确定合适的滞后期数。滞后期数的选择直接影响模型的预测能力和计算复杂性，滞后期数过多可能导致模型过拟合，而滞后期数过少可能无法捕捉到足够的动态信息。通常这些方法使用需要数据满足某些前提条件，但实际交通数据太复杂，无法满足这些前提条件，因此此类模型在实际运用中通常表现欠佳。


\subsection{机器学习预测方法}
由于经典统计方法不适合处理复杂和动态的时间序列数据，因此在实际应用中通常性能较差。越来越多的研究开始转向使用机器学习方法预测交通流量。机器学习方法在处理高维数据和揭示复杂的非线性关系时具有一定的优势。K近邻（KNN）和支持向量回归（SVR）是交通流量预测中常用的两种机器学习算法。

KNN模型在交通预测中的应用，关键在于如何选择合适的邻居相似度度量方法以及如何利用邻居的信息进行预测。\textcite{2012_于滨_K近邻短时交通流预测模型} 通过结合时间参数和空间参数，从时空两个维度衡量样本之间的距离。相比仅考虑时间因素的方法，该方式在预测性能上取得了显著提升。

SVM 通过核函数将数据映射到高维空间，从而捕捉到变量间复杂的非线性关系。基于SVR技术，\textcite{2013_Jeong_Supervised_Weighting_Online_Learning_Algorithm_for_Short_Term_Traffic_Flow_Prediction} 将权重参数的动态调整与支持向量回归法融合用于短期交通流量预测，权重参数的动态调整使得不同时段交通流特征差异性被充分考虑于该模型中。\textcite{2013_Lippi_Short_term_traffic_flow_forecasting_An_experimental_comparison_of_time_series_analysis_and_supervised_learning} 在SVR模型中通过添加描述季节性信息的函数，构建新RBF函数，辅助SVR模型学习交通的季节性信息，有效提高了预测精确性。

相较于经典统计模型，尽管机器学习方法在预测精度和使用灵活性上具有很大优势，但也存在一些显著的缺点：
\begin{enumerate}[(1)]
  \item 数据数量和质量上的强依赖性。充分的历史交通、天气、道路状况、节假日信息等信息在交通流量预测不可或缺。若数据稀缺或存在缺失，模型的预测能力将大大降低。例如，交通流量在某些地区或特定时间段的数据可能比较稀缺，这会导致模型无法有效地学习到数据的规律，影响其预测效果。此外，数据的质量至关重要，错误标注或噪声干扰会对模型的训练和预测产生较大的负面影响。
  \item 过度依赖人工特征选取，难以捕捉突发事件的隐含模式。
  \item 计算复杂度过高，计算资源和计算时间的开销过大。
  \item 模型泛化能力受限于训练数据覆盖面，易对局部数据波动产生过度反应，导致预测效果大打折扣。
\end{enumerate}


\subsection{深度学习预测方法}
在语音识别任务以及图像处理领域方面，深度学习技术已获得诸多突破性成果，由此引发研究者们对时空数据预测问题的关注度提升。具有强大表征学习能力的深度神经网络模型，正被越来越多地尝试运用于时空序列的建模过程当中，尤其是在交通流量、气象预报、能源需求、环境监测等领域取得了显著进展。时空数据预测不仅涉及时间维度的变化，还要考虑空间上的相互依赖性，因此深度学习为其提供了强大的工具。通过利用深度神经网络的高效特征学习能力，研究者能够从复杂的时空数据中提取出有效的特征模式，从而提高预测精度。

鉴于城市交通数据时间和空间上的相依性，卷积神经网络（CNN）、循环神经网络(RNN)、长短记忆神经网络(LSTM)被广泛使用 \parencite{2017_Zhang_Deep_spatio_temporal_residual_networks_for_citywide_crowd_flows_prediction, 1997_Hochreiter_Long_short_term_memory}。进一步，\textcite{2016_Wu_Shortterm_traffic_flow_forecasting_with_spatial_temporal_correlation_in_a_hybrid_deep_learning_framework} 及 \textcite{2018_Yao_Deep_multi_view_spatial_temporal_network_for_taxi_demand_prediction} 等研究提出融合CNN空间特征提取与LSTM时序建模能力的CLTFP框架。\textcite{2015_Shi_Convolutional_lstm_network_A_machine_learning_approach_for_precipitation_nowcasting} 提出一个具有嵌入卷积层的扩展全连接LSTM（FC-LSTM)。这些研究从不同方式融合卷积神经网络（CNN）和长短期记忆网络（LSTM），能够同时捕捉到空间上的局部特征和时间上的全局趋势，从而提升了交通流量预测的精度和鲁棒性。\textcite{2018_Yao_Modeling_spatial_temporal_dynamics_for_traffic_prediction} 提出了一种时空动态网络（ST-DN），用于出租车需求预测任务。该模型可以动态地学习不同地理位置之间的相似性和依赖关系，并根据实时交通和环境变化调整预测。

另一方面，为缓解深度网络中梯度消失问题，增强模型长期依赖的捕捉能力，\textcite{2018_Zhang_Predicting_citywide_crowd_flows_using_deep_spatio_temporal_residual_networks} 通过引入残差连接提出ST-ResNet 模型。

尽管这些深度学习模型在时空数据预测中取得了显著的成果，但它们仍然存在一定的局限性。首要限制在于多数模型对输入数据的格式要求严苛，需为标准的二维或三维网格结构。这种局限性导致模型难以直接处理具有非规则空间分布特性的城市交通流量数据。需将城市空间划分为固定尺寸的网格单元（如1km $\times$ 1km），这种人为的空间划分方式既无法适应离散化、异质化、动态变化的交通需求。例如，商业中心与远郊住宅区虽在地理位置上相距甚远，却可能共享相似的通勤高峰特征，这种跨区域的相关性超出了传统卷积神经网络（CNN）的建模能力范围。

在时间序列处理方面，现有方法主要依赖RNN及其衍生架构来捕获时序依赖关系。然而，这种链式结构在设计上存在固有缺陷：每个时间步的处理必须依次迭代，导致长序列建模时不可避免地出现信息衰减问题。此外，当前模型普遍缺乏对时间维度动态特性的深入考量，未能充分认识到前序时间步对当前预测的影响权重随时间变化的本质特征。这种静态的时间相关性处理方式，使得模型难以适应交通流量在不同时间段（如工作日与节假日、高峰与平峰时段）的动态演化规律，进一步制约了CNN+RNN混合架构在实际应用中的表现。这些技术瓶颈共同制约了深度学习模型在复杂交通场景中的泛化能力与鲁棒性，亟需发展新型时空建模框架以突破现有局限。

为了克服这些问题，近年来，研究者们开始探索更为灵活的深度学习架构，如图神经网络（GNN）和图卷积网络（GCN）。这些模型能够处理图形结构数据，在交通流量预测中尤为有效，因为城市交通网络本质上就是一个由路网、交通节点和路段组成的图形结构。通过引入图卷积网络，模型能够在非规则的交通网络上学习到时空依赖关系，从而实现更为精确的数据预测。基于图卷积网络（GCN）的短期交通流量预测也已取得令人瞩目的成就。然而，中长期（一小时及以上）的交通预测仍缺乏令人满意的进展，究其原因是存在以下不足和挑战：

\begin{enumerate}
  \item 动态时空建模局限性。
  道路网络中传感器之间的交通条件的相关性随着时间的推移而显著变化（例如，在高峰时间之前和期间）。这种相关性的动态变化导致基于静态图结构的方法难以准确建模交通流的真实演变过程。静态图无法及时反映节点间关系的时变特性，进而影响模型对复杂交通状态的感知与预测能力。因此，构建能够捕捉时空依赖变化的图模型，成为提升交通预测精度的关键方向。
  
  \item 对于误差传播的敏感性。
  当需要在一个长期尺度下进行预测时，人们主要使用的主体框架即为循环神经网络与图网络结合的方法，在该模型中每个时间步中的小误差可能会放大。这种错误的传播使预测对更长未来的时间片段具有高度的挑战性。
\end{enumerate}

近年来的研究中，注意机制被广泛应用于自然语言处理、图像字幕和语音识别等各种任务中。其核心目标在于，从所有输入信息中筛选出对于当前任务而言具有相对重要性的部分。在图像描述任务的研究过程中，两种不同类型的注意机制被 \textcite{2015_Xu_Show_attend_and_tell_Neural_image_caption_generation_with_visual_attention} 提出，同时采用可视化的技术手段直观呈现了该机制的运作效果。处理图结构数据时，基于自我注意层的神经网络架构被 \textcite{2018_Velickovic_Graph_attention_networks} 所采用，实例表明该方法在图节点分类任务中取得了突破性的成果。针对时间序列预测问题，一个包含多层次注意力结构的网络模型由 \textcite{2018_Liang_GeoMAN_Multi_level_Attention_Networks_for_Geo_sensory_Time_Series_Prediction} 设计完成，由此可见注意力机制能够自适应地学习多个地理传感器时序数据间的关联特性。

作为一种新兴技术，图注意力机制正逐步成为交通流量预测领域的关键技术。
交通流序列包括流量、速度、占有率等多个因素，它们之间相互作用，要准确预测交通流量，必须深入挖掘这些因素之间的非线性关系。\textcite{2019_Guo_Attention_based_spatial_temporal_graph_convolutional_networks_for_traffic_flow_forecasting}创新性地构建了一种基于注意力机制的新型时空图卷积网络模型（ASTGCN），该模型面向交通流量预测问题域展开探索。由三个相互独立的子模块构成ASTGCN模型的核心架构，这些子模块分别对应于交通流量数据中存在的三种典型时间模式特性：短期动态性、日周期性以及周周期性。具体而言，每个子模块内部又可划分为两个关键组成部分： 
\begin{enumerate}[(1)]
  \item 时空注意机制，以有效地捕获交通数据中的动态时空相关性。具体来说，空间注意被用来模拟不同位置之间复杂的空间相关性。采用时间注意法来捕捉不同时间之间的动态时间相关性；\item 时空卷积，通过时间卷积与空间卷积的双轨建模方式。它包括用于从原始的基于图的交通网络结构中捕获空间特征的图卷积，以及用于描述来自附近的时间切片的依赖关系的时间维数上的卷积。传统卷积运算的运用在此研究中呈现出新的可能性，与图神经网络形成互补态势。
\end{enumerate}

\textcite{2020_Zheng_GMAN_A_graph_multi_attention_network_for_traffic_prediction} 针对时空因素，提出了一种图注意网络（GMAN）来预测道路网络图上不同位置前方时间步长的交通情况。GMAN采用了一种编码器-解码器架构，其中时空注意模块构成编码器和解码器的主要组件。时空因素对交通状态的影响通过这种结构得以建模。流量特征编码由编码器完成，而输出序列预测任务则交由解码器执行。转换注意力层作为连接编码器与解码器的纽带，其功能是将已编码流量特征转化为未来时段的序列表征。作为输入解码器前的表征，历史时段与未来时段的直接关联通过转换注意机制建立。误差在预测时段间传播的现象由此得到有效抑制。在两个真实世界的交通预测任务（即交通量预测和交通速度预测）上的实验结果表明了GMAN的优越性。他们同样提出了空间和时间注意机制来建模动态空间和非线性时间相关性。此外，还设计了一种门控融合来自适应地融合由空间和时间注意机制提取的信息。

注意力机制对于图网络模型具有两方面的意义。传统的图卷积网络（GCN）在交通流量预测中常常会遇到“过平滑”问题，即多次图卷积操作后，节点的特征逐渐趋于一致，甚至相同。这种现象在交通流量预测中尤为不利，因为每个节点（如每个路段或区域）通常具有独特的时空特征。注意力机制的引入有利于解决这个问题。通过注意力机制，图神经网络能够根据节点之间的关系自适应地调整信息传播的权重，而不是简单地将所有邻接节点的特征平均或叠加。这样，模型能够在信息传播过程中，更加关注重要邻接节点，避免信息过度平滑和丢失。此外，注意力机制能够帮助网络识别出不同节点之间的不同影响程度。例如，在交通流量预测中，某些路段可能因交通管制、事故或突发事件而对预测结果产生更大影响，而其他路段则可能影响较小。注意力机制让模型自动学习到这些时空依赖关系，从而增强了模型对复杂交通网络的适应能力，提升了预测的准确性。

因此，结合注意力机制后，图神经网络能够在保持节点特征多样性的同时，有效避免过平滑问题，从而提升交通流量预测模型的性能。基于上述启发，本文主要的预测方法也将聚焦于图卷积与注意力机制的融合。

随着图网络研究的深入，一些研究开始对图卷积框架提出质疑。\textcite{2019_Wu_Graph_wavenet_for_deep_spatial_temporal_graph_modeling} 认为传统的预定义图模型未能充分捕捉复杂交通流预测场景中的空间依赖性，因此提出通过节点嵌入的方式构建自适应图邻接矩阵，为交通流预测中的空间关系建模提供了一种新思路。\textcite{2020_Bai_Adaptive_graph_convolutional_recurrent_network_for_traffic_forecasting} 则设计了节点自适应参数学习模块和数据自适应图生成模块，并与GRU网络结合，能够自动学习交通流数据之间的时空相关性。这两类方法的出现使得研究者可以跳出深奥的图理论，摆脱图的连通性信息的限制，也不再需要进行一些图拉普拉斯矩阵的运算，极大降低运算成本。本文将基于此思想改进融合图卷积与注意力机制框架。



\section{创新点}

本文创新性地提出了一种融合时空注意力机制与改进图卷积网络（GCN）的方法，以精准预测交通流量。该方法的主要创新之处体现在以下几个关键方面：

首先，本文引入了时间注意力机制，以弥补传统RNN和LSTM在预测序列中的不足。这些传统模型通常依赖固定的时间步长，难以捕捉交通流量在不同时间段的动态变化。相比之下，时间注意力机制通过自适应地为不同时间点分配权重，使模型能够聚焦于交通流量波动较大的时段，从而更精确地捕捉交通流量的时间变化规律。这一机制不仅提升了模型的预测准确性，还增强了对非线性和长时依赖关系的建模能力。

其次，本文采用空间注意力机制来缓解图卷积网络中的过平滑问题。在处理大规模交通图时，图卷积网络可能会面临不同节点表示过于相似的问题，导致模型失去辨别能力。为了解决这一问题，本研究引入了空间注意力机制，通过为图中的每个节点分配不同的注意力权重，使得节点间的信息传播更加灵活，有效避免了过度平滑现象的发生，进一步提升了模型的表现力。此外，在时空注意力机制与图网络的框架上进行了多信息的融合操作，验证了该框架对于多源信息融合的适应性。

在图网络设计方面，对图卷积操作进行了改进。考虑到交通网络规模庞大带来的计算开销，本文使用了一种先进的图滤波器结构，显著降低了计算复杂度。这种改进尤其适用于大规模图数据，有效减少了内存和计算资源的消耗，提高了计算效率。同时，该设计使得模型能够在更广泛的交通网络中进行有效预测。

此外，在论文的第二部分研究中提出了一种自适应图卷积操作。该操作摒弃了传统的预先定义的邻接矩阵，转而引入自适应学习机制，使模型能够根据实际数据中的时空依赖关系动态学习最合适的邻接矩阵。这一创新使得模型在进行图卷积时能够根据不同任务需求自动调整邻接矩阵，从而提高了模型的灵活性和表达能力。同时，由于不再依赖固定的邻接矩阵，计算量得以显著减少，进而加快了训练和推理过程，进一步提升了计算效率。



\section{本文结构}
现就本文的行文思路作如下总结： 

第一章为绪论部分。首先介绍交通流量预测的背景意义和应用前景；接着介绍交通流预测的国内国外研究现状；随后介绍论文的两个主要研究内容，主要的创新点和使用的技术手段；最后概述全文结构。 

第二章主要介绍相关模型的基本理论。章节首先介绍了图的相关理论支撑，图卷积背后的原理以及图卷积模块设计思路，为下文对于图网络层改进做铺垫。第二章还列举出交通流预测问题中的经典算法，并对其设计思路及原理进行简单介绍，作为后文基线模型使用。

第三章为本文的第一个主要部分。首先介绍基于时空注意力机制与图卷积融合的交通流预测方法整体流程框架；然后介绍网络的基本组成部分，阐明了每个部分的网络结构，将各个部分构建成为一个完整的交通流预测系统；最后，针对本章所提算法，在经典的交通流数据集上进行验证，将实验结果和现有方法进行比较，证明了该算法与之前的GCN+RNN框架相比，有效的提升了预测结果的准确度。随后，并提出了基于该框架下的多信息融合思路。

第四章为本文的第二个主要部分。章节首先介绍自适应图网络模型的交通流预测方法整体流程框架；然后对设计的动态图生成展开详细介绍；最后通过在多个数据集上进行对比实验，验证了算法在添加上述模块后有效提升了交通流数据的预测精度并减少推理时间。 

第五章对全文进行总结，并对所提方法的改进空间和未来研究方向提出展望。